\documentclass[11pt]{article}
\usepackage[letterpaper,margin=0.72in]{geometry}
\usepackage[T1]{fontenc}
\usepackage[utf8]{inputenc}
\usepackage{lmodern}
\usepackage{microtype}
\usepackage{xcolor}
\usepackage{graphicx}
\usepackage{booktabs}
\usepackage{longtable}
\usepackage{tabularx}
\usepackage{array}
\usepackage{enumitem}
\usepackage[most]{tcolorbox}
\usepackage{fancyhdr}
\usepackage{titlesec}
\usepackage{hyperref}
\usepackage{xurl}
\usepackage{lastpage}
\usepackage{ragged2e}
\usepackage{multicol}
\usepackage{amssymb}
\definecolor{BAOrange}{HTML}{E87521}
\definecolor{BANavy}{HTML}{17324D}
\definecolor{BASlate}{HTML}{5B6573}
\definecolor{BALight}{HTML}{F4F6F8}
\definecolor{BAGold}{HTML}{D7A62A}
\hypersetup{colorlinks=true,linkcolor=BANavy,urlcolor=BAOrange,citecolor=BANavy,pdftitle={Business Analysis Technical Guide}}
\setlength{\parindent}{0pt}
\setlength{\parskip}{4pt}
\setlist{nosep,leftmargin=*}
\renewcommand{\arraystretch}{1.18}
\newcolumntype{Y}{>{\RaggedRight\arraybackslash}X}
\newcolumntype{L}[1]{>{\RaggedRight\arraybackslash}p{#1}}
\pagestyle{fancy}
\setlength{\headheight}{14pt}
\fancyhf{}
\lhead{\textcolor{BANavy}{\small BABOK v3 Study Series}}
\rhead{\textcolor{BASlate}{\small Business Analysis Technical Guide}}
\cfoot{\small \thepage\ / \pageref{LastPage}}
\renewcommand{\headrulewidth}{0.4pt}
\titleformat{\section}{\Large\bfseries\color{BANavy}}{\thesection}{0.7em}{}
\titleformat{\subsection}{\large\bfseries\color{BAOrange}}{\thesubsection}{0.6em}{}
\titleformat{\subsubsection}{\normalsize\bfseries\color{BANavy}}{\thesubsubsection}{0.5em}{}
\newtcolorbox{keybox}[1][]{enhanced,breakable,colback=BALight,colframe=BANavy,boxrule=0.7pt,arc=1.5mm,left=2mm,right=2mm,top=1.5mm,bottom=1.5mm,#1}
\newtcolorbox{practicebox}[1][]{enhanced,breakable,colback=white,colframe=BAOrange,boxrule=0.8pt,arc=1.5mm,left=2mm,right=2mm,top=1.5mm,bottom=1.5mm,#1}
\newtcolorbox{warningbox}[1][]{enhanced,breakable,colback=white,colframe=BAGold,boxrule=0.8pt,arc=1.5mm,left=2mm,right=2mm,top=1.5mm,bottom=1.5mm,#1}
\begin{document}
\begin{titlepage}
\centering
\vspace*{1.15in}
{\Huge\bfseries\color{BANavy} Business Analysis Technical Guide\par}
\vspace{0.35in}
{\Large\color{BAOrange} Concepts, lifecycle, governance, value, and practical application\par}
\vspace{0.55in}
\rule{0.78\textwidth}{1.2pt}\par
\vspace{0.35in}
{\large Technical reference and study artifact based on the uploaded \emph{A Guide to the Business Analysis Body of Knowledge (BABOK Guide), Version 3}.\par}
\vfill
\begin{keybox}
\textbf{Source discipline.} BABOK terminology and organization are preserved where practical. Explanations are paraphrased for study and engineering use. Any engineering, security, Agile, UML/BPMN, or enterprise-architecture crosswalks are practitioner mappings, not normative BABOK statements.
\end{keybox}
\vspace{0.35in}
{\small Prepared October 2026\par}
\end{titlepage}
\tableofcontents
\newpage
\section{Purpose and Operating Model}
Business analysis is best understood as a value-oriented change discipline rather than a documentation function. The core job is to make a need explicit, understand the context and stakeholders around it, describe viable solutions, and help the enterprise decide and learn as evidence changes.

\begin{keybox}
\textbf{Working definition.} Business analysis enables change by clarifying needs and recommending solutions that deliver value to stakeholders. The guide applies this work before, during, and after projects and also to continuous enterprise evolution.
\end{keybox}

The framework is deliberately non-prescriptive. Knowledge areas are logical groupings of work, not lifecycle phases. Tasks can be performed sequentially, iteratively, or concurrently whenever the necessary information is available.

\subsection{The BABOK Structural Model}
\begin{itemize}
\item \textbf{Key concepts} establish the shared conceptual language.
\item \textbf{Knowledge areas} group 30 business analysis tasks into six domains of practice.
\item \textbf{Underlying competencies} describe the analytical, behavioral, business, communication, interaction, and tool skills that enable effective work.
\item \textbf{Techniques} provide reusable methods for performing tasks.
\item \textbf{Perspectives} adapt the same core work to common contexts such as Agile, IT, business architecture, business intelligence, and business process management.
\end{itemize}
\section{Business Analysis Core Concept Model}
The BACCM provides a compact system for checking whether analysis is complete enough to support a decision. None of the six concepts is sufficient by itself; each gains meaning through its relationship with the others.

\begin{longtable}{L{0.13\textwidth} L{0.37\textwidth} L{0.40\textwidth}}
\toprule
\textbf{Concept} & \textbf{Practical interpretation} & \textbf{Diagnostic questions}\\
\midrule
\endhead
Change & Deliberate transformation undertaken in response to a need. & What is changing? At what scope? What transition is required?\\
Need & A problem or opportunity that warrants attention. & What problem or opportunity exists? What evidence shows it matters?\\
Solution & A particular way of satisfying one or more needs in context. & What solution options exist? What assumptions make them viable?\\
Stakeholder & A person or group related to the need, change, or solution. & Who is affected, influential, accountable, or dependent?\\
Value & Worth, usefulness, benefit, loss, risk, or other consequence to a stakeholder. & What value is expected? For whom? How will we know it was realized?\\
Context & The circumstances that shape and are shaped by the change. & Which policies, culture, technology, processes, markets, constraints, and external forces matter?\\
\bottomrule
\end{longtable}

\begin{practicebox}
\textbf{BACCM review habit.} When a requirement, design decision, backlog item, or architecture proposal changes, re-check all six concepts. A local change can alter stakeholder impact, context, risk, solution feasibility, or expected value.
\end{practicebox}
\section{Requirements and Designs}
A requirement is a usable representation of a need; a design is a usable representation of a solution. The distinction is contextual: the same representation may function as a design for one stakeholder and as a requirement for another downstream stakeholder.

\subsection{Requirements Classification}
\begin{longtable}{L{0.20\textwidth} L{0.68\textwidth}}
\toprule
\textbf{Class} & \textbf{Meaning}\\
\midrule
\endhead
Business requirements & Goals, objectives, and outcomes explaining why the change is being initiated.\\
Stakeholder requirements & Stakeholder needs that must be satisfied to achieve the business requirements.\\
Solution requirements & Capabilities and qualities of the solution. These include functional requirements and non-functional/quality-of-service requirements.\\
Transition requirements & Temporary capabilities or conditions required to move from current to future state, such as conversion, training, migration, and business continuity.\\
\bottomrule
\end{longtable}

\begin{warningbox}
\textbf{Common failure mode.} Teams often begin with a proposed solution and immediately write implementation requirements. Preserve the chain of reasoning: business need $\rightarrow$ stakeholder need $\rightarrow$ solution requirement $\rightarrow$ design $\rightarrow$ evidence of value.
\end{warningbox}
\section{The Six Knowledge Areas}
The six knowledge areas form an interacting system. Planning establishes how analysis will work; elicitation and collaboration supply and reconcile information; strategy explains why and where to change; requirements life-cycle management governs information over time; analysis and design definition turn information into coherent requirements and options; solution evaluation closes the loop with evidence of realized value.

\subsection{Business Analysis Planning and Monitoring}
Organize and coordinate business analysis work, stakeholder participation, governance, information management, and improvement.
\begin{itemize}
\item \textbf{3.1 Plan Business Analysis Approach}: Define how business analysis will be performed for the initiative.
\item \textbf{3.2 Plan Stakeholder Engagement}: Plan how to establish and maintain productive stakeholder relationships.
\item \textbf{3.3 Plan Business Analysis Governance}: Define how requirements and design decisions, changes, priorities, reviews, and approvals will be governed.
\item \textbf{3.4 Plan Business Analysis Information Management}: Define how business analysis information will be captured, organized, stored, accessed, traced, and retained.
\item \textbf{3.5 Identify Business Analysis Performance Improvements}: Assess the effectiveness of business analysis work and identify improvements.
\end{itemize}
\subsection{Elicitation and Collaboration}
Obtain information, confirm what was learned, communicate it effectively, and sustain stakeholder collaboration throughout the work.
\begin{itemize}
\item \textbf{4.1 Prepare for Elicitation}: Set the scope of an elicitation activity, select techniques, and arrange the people, materials, and logistics needed.
\item \textbf{4.2 Conduct Elicitation}: Explore and draw out information relevant to the change.
\item \textbf{4.3 Confirm Elicitation Results}: Check elicited information for accuracy and consistency before relying on it.
\item \textbf{4.4 Communicate Business Analysis Information}: Present business analysis information so stakeholders develop a shared understanding.
\item \textbf{4.5 Manage Stakeholder Collaboration}: Sustain stakeholder participation and alignment toward common outcomes.
\end{itemize}
\subsection{Requirements Life Cycle Management}
Manage requirements and designs from inception through retirement, preserving alignment, traceability, priority, change control, and agreement.
\begin{itemize}
\item \textbf{5.1 Trace Requirements}: Maintain relationships among needs, requirements, designs, solution components, and related work so changes and coverage can be assessed.
\item \textbf{5.2 Maintain Requirements}: Keep requirements and designs accurate, current, consistent, and reusable across their life cycle.
\item \textbf{5.3 Prioritize Requirements}: Order requirements and designs according to relative importance, value, urgency, risk, and dependencies.
\item \textbf{5.4 Assess Requirements Changes}: Evaluate the consequences of proposed changes to requirements and designs.
\item \textbf{5.5 Approve Requirements}: Obtain the level of agreement and authorization needed to continue analysis or proceed with solution work.
\end{itemize}
\subsection{Strategy Analysis}
Understand the business need and current state, define the desired future state, assess uncertainty, and select a change strategy.
\begin{itemize}
\item \textbf{6.1 Analyze Current State}: Establish why change is needed and understand the existing enterprise conditions affected by it.
\item \textbf{6.2 Define Future State}: Describe the conditions, goals, objectives, and capabilities that would satisfy the business need.
\item \textbf{6.3 Assess Risks}: Evaluate uncertainty that could affect value during transition or in the future state, and identify suitable responses.
\item \textbf{6.4 Define Change Strategy}: Compare ways to move from current to future state and recommend the most valuable feasible approach, including transition states.
\end{itemize}
\subsection{Requirements Analysis and Design Definition}
Transform elicited information into coherent requirements and designs, assess quality and value, organize the architecture, and recommend a solution.
\begin{itemize}
\item \textbf{7.1 Specify and Model Requirements}: Analyze, synthesize, and refine elicitation results into useful requirements and design representations.
\item \textbf{7.2 Verify Requirements}: Check that requirements and designs meet quality expectations and are fit for their intended analytical or delivery purpose.
\item \textbf{7.3 Validate Requirements}: Confirm that requirements and designs support business requirements and the delivery of needed value.
\item \textbf{7.4 Define Requirements Architecture}: Organize requirements so their relationships and collective structure support the objectives.
\item \textbf{7.5 Define Design Options}: Develop solution approaches, allocate requirements, and represent viable options for achieving the future state.
\item \textbf{7.6 Analyze Potential Value and Recommend Solution}: Estimate the potential value of design options and recommend the option best aligned with enterprise needs.
\end{itemize}
\subsection{Solution Evaluation}
Measure how an implemented or partially implemented solution performs, diagnose limitations, and recommend actions that increase realized value.
\begin{itemize}
\item \textbf{8.1 Measure Solution Performance}: Define meaningful measures and collect data to evaluate how effectively the solution delivers value.
\item \textbf{8.2 Analyze Performance Measures}: Interpret performance information to understand realized value and whether current business needs are being met.
\item \textbf{8.3 Assess Solution Limitations}: Identify internal solution factors that constrain value realization.
\item \textbf{8.4 Assess Enterprise Limitations}: Identify external enterprise factors that prevent full value realization.
\item \textbf{8.5 Recommend Actions to Increase Solution Value}: Recommend actions that close the gap between potential and actual value.
\end{itemize}
\section{A Practical End-to-End Analysis Loop}
The following loop is not an official BABOK lifecycle; it is a practitioner way to combine the non-sequential tasks into a coherent working rhythm.

\begin{enumerate}
\item \textbf{Frame the change.} Establish the business need, current state, stakeholders, context, and desired outcomes.
\item \textbf{Plan the analysis system.} Decide methods, governance, information management, traceability, and stakeholder engagement.
\item \textbf{Elicit and model.} Obtain information using interviews, workshops, observation, document analysis, prototypes, and models appropriate to the problem.
\item \textbf{Specify and test the reasoning.} Structure requirements, verify their quality, validate their value alignment, and organize their relationships.
\item \textbf{Explore solution options.} Allocate requirements, compare designs, risks, costs, and potential value, then recommend an option.
\item \textbf{Govern change.} Maintain, prioritize, trace, assess changes, and obtain appropriate approvals as information evolves.
\item \textbf{Evaluate actual outcomes.} Measure the operating solution, identify internal and enterprise limitations, and recommend improvements.
\item \textbf{Learn and repeat.} Feed performance results and lessons into the next planning, strategy, and elicitation cycle.
\end{enumerate}
\section{Underlying Competencies}
\begin{longtable}{L{0.24\textwidth} L{0.64\textwidth}}
\toprule
\textbf{Competency group} & \textbf{Engineering interpretation}\\
\midrule
Analytical Thinking and Problem Solving & Decompose ambiguous problems, reason about cause and effect, evaluate alternatives, learn quickly, and apply systems thinking.\\
Behavioural Characteristics & Act ethically and accountably; build trust; remain adaptable; organize work and commitments reliably.\\
Business Knowledge & Understand the enterprise, industry, solution domain, methods, and business drivers well enough to judge value and feasibility.\\
Communication Skills & Listen accurately; write and speak clearly; tailor information to audience, context, and decision need.\\
Interaction Skills & Facilitate groups, negotiate, resolve conflict, lead without relying on authority, and support shared decisions.\\
Tools and Technology & Use modelling, requirements, collaboration, communication, productivity, and analytical tools appropriately.\\
\bottomrule
\end{longtable}
\section{Technique Selection}
A technique should be chosen because it produces the evidence, representation, or stakeholder interaction needed for the current task. Useful selection questions include:
\begin{itemize}
\item Is the problem primarily about \textbf{discovery}, \textbf{modelling}, \textbf{decision making}, \textbf{measurement}, \textbf{risk}, or \textbf{collaboration}?
\item How much uncertainty exists and how expensive is it to be wrong?
\item Which stakeholders must participate, and what is their availability and decision authority?
\item Is the information best expressed as text, a table, a process, a state machine, an interaction, a data model, or a prototype?
\item What level of traceability and governance is required?
\end{itemize}

The companion \emph{50 Techniques Catalog} provides a quick reference to every technique in Chapter 10.
\section{Perspectives}
Perspectives are lenses for tailoring business analysis to a context. They do not replace the six knowledge areas.
\begin{longtable}{L{0.20\textwidth} L{0.68\textwidth}}
\toprule
\textbf{Perspective} & \textbf{Primary emphasis}\\
\midrule
Agile & Continuous analysis, incremental delivery, prioritization, close team collaboration, and just-in-time detail.\\
Business Intelligence & Information, analytics, data quality, decision support, and value obtained from data.\\
Information Technology & Effects of change on information systems and translation of needs into technology-oriented solutions.\\
Business Architecture & Strategy, capabilities, operating model, enterprise scope, and alignment of change with organizational direction.\\
Business Process Management & Process-centric improvement of how work is performed and governed across functional boundaries.\\
\bottomrule
\end{longtable}
\section{Operational Checklists}
\subsection{Before analysis begins}
\begin{itemize}
\item Business need and decision to be supported are explicit.
\item Stakeholders, decision rights, and collaboration expectations are known.
\item Information management and traceability conventions are defined.
\item Current-state evidence is sufficient to avoid designing from assumptions alone.
\end{itemize}
\subsection{Before requirements are accepted for delivery}
\begin{itemize}
\item Requirement class and source need are clear.
\item Requirement is unambiguous enough to support the intended decision or implementation.
\item Relevant dependencies, interfaces, quality attributes, risks, and constraints are represented.
\item Acceptance/evaluation criteria are testable or otherwise measurable.
\item Requirement is traced to business/stakeholder value and to downstream design/test evidence as appropriate.
\end{itemize}
\subsection{After release or implementation}
\begin{itemize}
\item Performance measures are tied to the intended value, not only activity volume.
\item Solution and enterprise limitations are considered separately.
\item Improvement recommendations feed back into strategy, requirements, and prioritization.
\end{itemize}

\section*{Primary Source}
\addcontentsline{toc}{section}{Primary Source}
International Institute of Business Analysis (IIBA), \emph{A Guide to the Business Analysis Body of Knowledge (BABOK Guide)}, Version 3, 2015. This document is a study and practitioner synthesis of the user-supplied copy; it is not an official IIBA publication.
\end{document}
